% =====================================================================
%  Module 1 : Introduction à AWS
% =====================================================================
\section{Introduction à AWS}

\begin{frame}{Amazon Web Services}
  \begin{itemize}
    \setlength{\itemsep}{.55em}
    \item Filiale cloud d'Amazon, lancée en 2006
    \item Infrastructure et services informatiques à la demande, accessibles par API
    \item Premier fournisseur mondial de cloud : 28 \% du marché (T2 2026)
    \item Chiffre d'affaires 2025 : 128,7 milliards de dollars (+20 \% sur un an)
  \end{itemize}
  \bigskip
  \textbf{Principes fondateurs}
  \begin{itemize}
    \item API d'abord : tout se pilote par programme
    \item Libre-service : ressources disponibles en quelques minutes, sans intervention humaine
    \item Paiement à l'usage : à la seconde, au gigaoctet, à la requête
  \end{itemize}
  \source{Amazon, \emph{Fourth Quarter 2025 Results} ; Synergy Research Group, Q2 2026.}
\end{frame}

\begin{frame}{Quelques dates}
  \renewcommand{\arraystretch}{1.3}
  \begin{tabular}{@{}l l@{}}
    \textbf{2002} & premiers services web d'Amazon pour les développeurs \\
    \textbf{2006} & S3 (mars), SQS (juillet), EC2 (août) : les trois services d'origine \\
    \textbf{2009} & VPC et RDS \\
    \textbf{2012} & DynamoDB \\
    \textbf{2014} & Lambda (fonctions serverless) \\
    \textbf{2017} & ouverture de la région de Paris (\texttt{eu-west-3}) \\
    \textbf{2026} & 39 régions, 123 zones de disponibilité, plus de 1 200 types d'instances EC2 \\
  \end{tabular}
  \source{AWS News Blog, « Twenty years of Amazon S3 », « Happy 20th Birthday, Amazon EC2 », « Amazon SQS turns 20 », 2026.}
\end{frame}

\begin{frame}{Parts de marché du cloud (T2 2026)}
  \begin{columns}[c]
    \begin{column}{.55\textwidth}
      \begin{tikzpicture}[x=1cm,y=1cm]
        \foreach \nom/\part/\coul [count=\i from 0] in {AWS/28/Orange, Microsoft Azure/20/Lien, Google Cloud/15/Vert, Autres/37/Ligne}{
          \node[anchor=east, font=\small] at (0,-\i*.9) {\nom};
          \fill[\coul] (.15,-\i*.9-.28) rectangle ++(\part*.1,.56);
          \node[anchor=west, font=\small\bfseries] at (\part*.1+.25,-\i*.9) {\part\ \%};
        }
      \end{tikzpicture}
    \end{column}
    \begin{column}{.42\textwidth}
      \begin{itemize}
        \setlength{\itemsep}{.5em}
        \item Marché mondial de l'infrastructure cloud : 143,4 Md\$ sur le trimestre
        \item Croissance d'environ 43 \% sur un an
        \item Les trois premiers : 63 \% du marché
      \end{itemize}
    \end{column}
  \end{columns}
  \source{Synergy Research Group, « Cloud Market Share Trends », Q2 2026.}
\end{frame}

\begin{frame}{Modèles de service}
  \renewcommand{\arraystretch}{1.45}
  \begin{tabularx}{\textwidth}{@{}l>{\raggedright\arraybackslash}X>{\raggedright\arraybackslash}X@{}}
    \toprule
    & \textbf{Le client gère} & \textbf{Exemples AWS} \\ \midrule
    \textbf{IaaS} & système d'exploitation, logiciels, données, accès & EC2, EBS, VPC \\
    \textbf{PaaS} & code, données, accès & RDS, Lambda, Elastic Beanstalk \\
    \textbf{SaaS} & données, accès & Amazon WorkMail, Amazon Connect \\ \bottomrule
  \end{tabularx}
  \bigskip
  \begin{itemize}
    \item Plus le service est géré, moins il y a d'administration, et moins de contrôle
    \item Les données et les droits d'accès restent toujours à la charge du client
  \end{itemize}
\end{frame}

\begin{frame}{Modèle de responsabilité partagée}
  \begin{columns}[c]
    \begin{column}{.6\textwidth}
      \includegraphics[width=\textwidth]{responsabilite}
    \end{column}
    \begin{column}{.38\textwidth}
      \begin{itemize}
        \setlength{\itemsep}{.6em}
        \item AWS : sécurité \textbf{du} cloud (bâtiments, matériel, réseau, virtualisation)
        \item Client : sécurité \textbf{dans} le cloud (configuration, mises à jour, données, IAM)
        \item La frontière dépend du service
      \end{itemize}
    \end{column}
  \end{columns}
  \lien{https://aws.amazon.com/compliance/shared-responsibility-model/}
\end{frame}

\begin{frame}{Catalogue de services}
  \centering
  \includegraphics[height=.74\textheight]{familles-services}
  \source{AWS, \emph{Overview of Amazon Web Services} : plus de 200 services.}
\end{frame}

\begin{frame}{Infrastructure classique et équivalents AWS}
  \centering
  \includegraphics[width=.92\textwidth]{infra-correspondance}
\end{frame}

\begin{frame}{Infrastructure mondiale : régions}
  \begin{itemize}
    \setlength{\itemsep}{.55em}
    \item Région : zone géographique regroupant plusieurs zones de disponibilité
    \item 39 régions en 2026, d'autres annoncées
    \item Régions indépendantes : une panne reste en principe confinée à sa région
    \item Les données restent dans la région choisie, sauf réplication explicite
    \item Chaque région a un code : \texttt{eu-west-3} (Paris), \texttt{eu-central-1} (Francfort), \texttt{us-east-1} (Virginie du Nord)
    \item Certaines régions récentes doivent être activées (\emph{opt-in}) : \texttt{af-south-1}, \texttt{eu-south-1}...
  \end{itemize}
  \lien{https://aws.amazon.com/about-aws/global-infrastructure/regions_az/}
\end{frame}

\begin{frame}{Régions, zones de disponibilité, points de présence}
  \centering
  \includegraphics[width=\textwidth]{region-az}
\end{frame}

\begin{frame}{Zones de disponibilité (AZ)}
  \begin{itemize}
    \setlength{\itemsep}{.55em}
    \item Un ou plusieurs centres de données à l'intérieur d'une région
    \item Alimentation, refroidissement et réseau propres
    \item Distantes les unes des autres, mais à moins de 100 km : latence de l'ordre de la milliseconde
    \item Paris : 3 AZ, \texttt{eu-west-3a}, \texttt{eu-west-3b}, \texttt{eu-west-3c}
    \item Points de présence : plusieurs centaines de sites pour CloudFront et Route 53
  \end{itemize}
  \bigskip
  \attention Un sous-réseau, un volume EBS et une instance EC2 vivent dans \textbf{une seule} AZ
\end{frame}

\begin{frame}{Haute disponibilité}
  \centering
  \includegraphics[width=.9\textwidth]{haute-dispo}\par\bigskip
  \raggedright
  \begin{itemize}
    \item Haute disponibilité = ressources réparties sur au moins deux AZ
    \item Pour EC2 : plusieurs instances derrière un répartiteur de charge (ELB)
    \item Pour une base RDS : déploiement Multi-AZ (copie de secours synchrone)
  \end{itemize}
\end{frame}

\begin{frame}{Services et régions}
  \begin{itemize}
    \setlength{\itemsep}{.55em}
    \item Tous les services ne sont pas disponibles dans toutes les régions
    \item Les nouveautés arrivent souvent d'abord en \texttt{us-east-1}
    \item Les prix varient d'une région à l'autre
    \item Services \textbf{régionaux} : EC2, EBS, VPC, RDS, SQS... (la console affiche une région)
    \item Services \textbf{globaux} : IAM, Route 53, CloudFront (la console affiche « Global »)
    \item S3 : liste des buckets globale, mais chaque bucket appartient à une région
  \end{itemize}
  \bigskip
  \attention Une ressource créée dans une autre région n'apparaît pas dans la console
  \lien{https://aws.amazon.com/about-aws/global-infrastructure/regional-product-services/}
\end{frame}

\begin{frame}{Choisir une région}
  \begin{enumerate}
    \setlength{\itemsep}{.6em}
    \item \textbf{Réglementation} : localisation des données (RGPD, données de santé, secteur public)
    \item \textbf{Latence} : proximité des utilisateurs, à mesurer
    \item \textbf{Services disponibles} dans la région
    \item \textbf{Prix}
  \end{enumerate}
  \bigskip
  Pour ce cours : \texttt{eu-west-3} (Paris), ouverte en décembre 2017, 3 AZ
\end{frame}

\begin{frame}{Idée reçue n°1 : « le cloud n'est pas fiable »}
  \begin{itemize}
    \setlength{\itemsep}{.5em}
    \item Engagements de service (SLA) publiés pour chaque service
    \begin{itemize}
      \item EC2 : 99,99 \% par mois pour un déploiement sur plusieurs AZ, 99,5 \% pour une instance seule
      \item S3 Standard : 99,9 \% de disponibilité, durabilité conçue pour 99,999999999 \% (11 neuf)
    \end{itemize}
    \item Données S3 copiées dans plusieurs AZ
  \end{itemize}
  \bigskip
  \attention Les pannes existent :
  \begin{itemize}
    \item 28/02/2017 : S3 indisponible plus de 4 h en \texttt{us-east-1} (erreur de commande)
    \item 19-20/10/2025 : panne de DynamoDB d'environ 14 h en \texttt{us-east-1}, avec effet en cascade sur EC2, Lambda, STS...
  \end{itemize}
  \source{AWS, \emph{Amazon Compute SLA}, \emph{Amazon S3 SLA} ; AWS, résumés des incidents du 28/02/2017 et du 19-20/10/2025.}
\end{frame}

\begin{frame}{Idée reçue n°2 : « je perds le contrôle de mes données »}
  \begin{itemize}
    \setlength{\itemsep}{.6em}
    \item Le client reste propriétaire de ses données
    \item Il choisit la région où elles sont stockées
    \item Chiffrement au repos et en transit, avec des clés gérées par AWS ou par le client (KMS)
    \item Récupération des données à tout moment (API, transfert, AWS Snowball)
    \item Journal de tous les appels d'API du compte (CloudTrail)
  \end{itemize}
\end{frame}

\begin{frame}{Tarification}
  \begin{itemize}
    \setlength{\itemsep}{.5em}
    \item Paiement à l'usage, facture mensuelle
    \item Chaque service a ses propres unités : heure ou seconde, Go-mois, requêtes, Go transférés
    \item Coût difficile à prévoir sans estimation préalable
    \item Pas d'arrêt automatique en cas de dépassement : un budget AWS envoie une alerte, il ne coupe rien
    \item Comptes créés depuis le 15/07/2025 : jusqu'à 200 \$ de crédits, plan gratuit de 6 mois
    \item Transfert entrant gratuit, transfert sortant vers Internet facturé
  \end{itemize}
  \lien{https://calculator.aws}
\end{frame}

\begin{frame}{Ce qui coûte, même sans être utilisé}
  \renewcommand{\arraystretch}{1.3}
  \begin{tabularx}{\textwidth}{@{}>{\raggedright\arraybackslash}p{4.4cm}>{\raggedright\arraybackslash}X>{\raggedright\arraybackslash}p{3.8cm}@{}}
    \toprule
    \textbf{Ressource} & \textbf{Unité facturée} & \textbf{Facturée à l'arrêt ?} \\ \midrule
    Instance EC2 démarrée & seconde (minimum 60 s) & \textcolor{Rouge}{oui} \\
    Instance EC2 arrêtée & aucune pour le calcul & \textcolor{Rouge}{disque EBS facturé} \\
    Volume EBS & Go provisionné par mois & \textcolor{Rouge}{oui} \\
    Adresse IPv4 publique & heure (0,005 \$) depuis 02/2024 & \textcolor{Rouge}{oui} \\
    NAT Gateway & heure + Go traité & \textcolor{Rouge}{oui} \\
    Objet S3 & Go-mois + requêtes & selon le volume \\ \bottomrule
  \end{tabularx}
  \source{AWS, pages de tarification ; AWS News Blog, « New – AWS Public IPv4 Address Charge », 2023.}
\end{frame}

\begin{frame}{Accès à AWS}
  \centering
  \includegraphics[height=.56\textheight]{appel-api}\par\medskip
  \raggedright
  \begin{itemize}
    \item Console web : découverte, diagnostic
    \item CLI (\texttt{aws}) et SDK (\texttt{boto3}...) : scripts et applications
    \item Infrastructure as Code : CloudFormation, CDK, Terraform, Pulumi
  \end{itemize}
\end{frame}

\begin{frame}[fragile]{AWS CLI et CloudShell}
  \begin{itemize}
    \setlength{\itemsep}{.5em}
    \item CloudShell : terminal intégré à la console, déjà authentifié, CLI installée
    \item Toute action passe par un appel d'API signé, contrôlé par IAM
    \item Même appel, que l'on passe par la console, la CLI ou un SDK
  \end{itemize}
  \begin{terminal}
aws sts get-caller-identity\\
aws ec2 describe-availability-zones --region eu-west-3 --output table
  \end{terminal}
  \lien{https://docs.aws.amazon.com/cli/latest/reference/}
\end{frame}

\diapotp{1}{Prise en main de la console}{%
  \item Connexion avec l'utilisateur IAM fourni
  \item Région par défaut : \texttt{eu-west-3}
  \item Premières commandes dans CloudShell
  \item Lecture des prix des types d'instances}

